% ============================================================
% PHẦN B - GIỚI THIỆU SCIMKG
% ============================================================
\section{Giới thiệu}\label{sec:sm-intro}

Học liệu mở như khóa học trực tuyến đại trà (MOOC), slide và video bài giảng tăng rất nhanh, nhưng ở dạng rời rạc nên khó dùng lại. Đồ thị tri thức giáo dục (EKG) gom nội dung đó thành các khái niệm và quan hệ giữa chúng, lưu dưới dạng triple (đầu, quan hệ, đuôi). Nhờ cấu trúc này, các ứng dụng như sinh câu hỏi, hỏi đáp, truy hồi tri thức hay trợ lý học tập có chỗ dựa để tra cứu.

Theo \SM{}~\cite{lu2026scimkg}\footnote{\url{https://doi.org/10.1609/aaai.v40i18.38574}}, các EKG hiện có còn hai hạn chế. Thứ nhất, nhiều hệ thống phải gán nhãn dữ liệu (hoặc giám sát yếu) và huấn luyện mô hình riêng để rút khái niệm, như KnowEdu~\cite{chen2018knowedu} và MOOCCubeX~\cite{yu2021mooccubex}, nên khó theo kịp nguồn học liệu liên tục thay đổi. Ngay cả các cách dùng LLM gần đây, như của Jhajj và cộng sự~\cite{jhajj2024gptekg} và LLM4EduKG~\cite{sun2024llm4edukg}, theo tác giả vẫn phụ thuộc một phần vào khung đồ thị dựng tay. Thứ hai, phần lớn đồ thị chỉ chứa văn bản, còn tri thức nằm trong ảnh, video hay lời nói bị bỏ qua hoặc quy hết về chữ. TIVA-KG~\cite{wang2023tivakg} có đủ bốn phương thức, nhưng căn chỉnh thủ công và thuộc miền tổng quát.

Để giải quyết hai hạn chế này, tác giả đề xuất dựng EKG đa phương thức gần như tự động, chỉ bằng cách gọi các mô hình ngôn ngữ lớn (LLM) và mô hình ngôn ngữ đa phương thức (MLLM) có sẵn, không huấn luyện hay tinh chỉnh trọng số. Chuỗi EVIA (Extraction, Verification, Integration, Augmentation) cho nhiều LLM cùng đọc từng đoạn bài giảng, tự phản hồi và bỏ phiếu để chỉ giữ những khái niệm thực sự thuộc môn học. Mỗi khái niệm sau đó được gắn với ảnh, đoạn video, lời giải thích và âm thanh tương ứng. Trên ba môn sinh học, vật lý và hóa học bậc phổ thông của Trung Quốc, đồ thị thu được có 1.356 điểm kiến thức (knowledge point), 34.630 khái niệm và 403.400 triple. Về trích xuất khái niệm, bài báo công bố F1 đạt 0,803, so với 0,734 của baseline tốt nhất (LinkNER). Theo tác giả, đây là EKG đầu tiên bao phủ đồng thời văn bản, ảnh, video và âm thanh.

Báo cáo mô tả \SM{} theo đúng thứ tự dữ liệu đi qua hệ thống: dữ liệu vào, mô hình ở mỗi bước, vai trò của từng prompt và cách tính điểm. Những chỗ bài báo chỉ nêu ở mức khái quát được bổ sung từ mã nguồn công bố\footnote{\url{https://github.com/kg-bnu/SciMKG}}, và nhóm nêu rõ các điểm bài báo và mã không khớp. Tiếp theo, báo cáo tổng hợp kết quả công bố và đánh giá giới hạn của cả phương pháp lẫn cách đánh giá.

Phần còn lại của Phần B gồm công trình liên quan, phương pháp, thực nghiệm và kết luận.
